\chapter{读懂数学证明的入门语法}

学一门语言，先要学会它的语法；读数学证明也一样。证明并不是一堆符号的杂技，而是一套“把话说清楚”的规则：什么叫一句话成立，什么叫一句话被推翻，什么样的理由才算数。本附录把全书反复使用的这套语法集中在一起，配上初中水平的例子。读正文时如果哪里卡住了，随时可以翻回来查。

\section{命题：可以判断真假的话}

数学只讨论一种特殊的句子：\textbf{命题}，也就是能够明确判断对或错的陈述。“三角形的内角和是 $180^\circ$”是命题（它是对的）；“所有的质数都是奇数”也是命题（它是错的，$2$ 就是反例）。但“数学真有意思”不是命题，因为没法判定真假；“$x>3$”单独看也不是命题，因为 $x$ 没确定时无法判断——这正是量词要解决的问题（见后文）。

很多命题有“如果……那么……”的形状。比如“如果一个整数能被 $4$ 整除，那么它能被 $2$ 整除”。“如果”后面的部分叫条件，“那么”后面的部分叫结论。把条件和结论对调，得到的新句子叫原命题的\textbf{逆命题}。上面那句话的逆命题是：“如果一个整数能被 $2$ 整除，那么它能被 $4$ 整除。”请注意：原命题是对的，逆命题却是错的（$6$ 能被 $2$ 整除，却不能被 $4$ 整除）。

\begin{fanli}
“原命题成立，逆命题就成立”是最常见的误会之一。“对顶角相等”是对的，但它的逆命题“相等的角是对顶角”明显错误——两直线平行时的同位角相等，却根本不是对顶角。判断一句话对不对，和判断它的逆命题对不对，是两笔独立的账，必须分开算。
\end{fanli}

\section{充分条件与必要条件：两顶不能混戴的帽子}

当“如果 $p$，那么 $q$”成立时，我们说 $p$ 是 $q$ 的\textbf{充分条件}：有 $p$ 就“足够”保证 $q$。同时说 $q$ 是 $p$ 的\textbf{必要条件}：想要 $p$ 成立，$q$ 是“必须”跨过的门槛。

\begin{li}
“四边形是矩形”能保证“它是平行四边形”，所以“是矩形”是“是平行四边形”的充分条件；反过来，“是平行四边形”只是“是矩形”的必要条件——矩形首先得是平行四边形，但光平行还不够，还必须有直角。
\end{li}

一个生活类比：有电是灯泡发光的必要条件（没电一定不亮），但不是充分条件（有电，灯丝断了也不亮）。充分不一定必要，必要不一定充分；两者都满足时，我们称 $p$ 是 $q$ 的\textbf{充要条件}，记作 $p\Leftrightarrow q$，这时两者“同生同灭”，比如“三角形是等边三角形”与“它的三个内角都是 $60^\circ$”。

\section{量词：把“所有”和“存在”说清楚}

命题里最常见的两个词是“所有”（任意，记作 $\forall$）和“存在”（至少有一个，记作 $\exists$）。它们的行为很不一样：

\begin{itemize}
\item “所有偶数都能被 $2$ 整除”要对，必须\textbf{每一个}偶数都过关；
\item “存在整数 $x$ 使 $x^2=4$”要对，只要\textbf{找到}一个就行（$x=2$ 即可）；
\item 推翻“所有”型命题，只需要一个反例；推翻“存在”型命题，却需要论证“一个都找不到”。
\end{itemize}

否定它们时规则也很干脆：“所有乌鸦都是黑的”的否定不是“所有乌鸦都不黑”，而是“存在一只乌鸦不是黑的”。全称的否定是特称，特称的否定是全称——这个交换律在反证法里天天要用。

\section{反例：一个就足够}

“反例”是满足条件却不满足结论的具体例子。它是数学里最便宜的武器：再漂亮的猜想，一个反例就能推翻。猜想“对所有整数 $a,b$，若 $a^2>b^2$ 则 $a>b$”，听上去挺像回事，但取 $a=-3$，$b=1$：$a^2=9>1=b^2$，可是 $-3<1$。一个反例，猜想作废。反过来，验证了一百万个例子也不能算证明——例子只能说明“还没被打倒”，不能保证“永远不会被打倒”。

\section{数学归纳法：推倒一排多米诺骨牌}

想证明一个关于所有正整数 $n$ 的命题 $P(n)$，比如
\[1+3+5+\cdots+(2n-1)=n^2,\]
一个个验算永远验不完。数学归纳法的想法是：只需做两件事——（1）\textbf{起点}：验证 $P(1)$ 成立；（2）\textbf{传递}：证明只要 $P(k)$ 成立，就能推出 $P(k+1)$ 成立。这就像推倒一排多米诺骨牌：第一块倒了，且每块倒下都会碰倒下一块，那么整排必倒。

\begin{proof}
（起点）$n=1$ 时，左边 $=1$，右边 $=1^2=1$，成立。
（传递）假设 $n=k$ 时 $1+3+\cdots+(2k-1)=k^2$ 成立，则 $n=k+1$ 时
\[1+3+\cdots+(2k-1)+(2k+1)=k^2+(2k+1)=(k+1)^2,\]
正好是要证的式子。由归纳法，等式对所有正整数 $n$ 成立。
\end{proof}

注意两步缺一不可：只验证起点，可能半路断掉；只证明传递，可能整排根本没被推倒。

\section{反证法：先假设敌人是对的}

有些命题正面进攻很难，那就换个打法：\textbf{假设结论不成立，往下推理，直到撞上矛盾}。矛盾说明开头的假设错了，于是结论只能成立。经典例子是“质数有无穷多个”。

\begin{proof}
假设质数只有有限个，把它们全部列出来：$p_1,p_2,\dots,p_n$。造一个数 $N=p_1p_2\cdots p_n+1$。$N$ 除以任何一个 $p_i$ 都余 $1$，所以没有一个已知质数能整除 $N$。但任何大于 $1$ 的整数要么本身是质数，要么能被某个质数整除，于是必然存在一个不在清单上的质数——这与“清单已列出全部质数”矛盾。所以质数不可能只有有限个。
\end{proof}

早在两千多年前，欧几里得的《几何原本》里就有这个论证。反证法的威力在于：它允许你“借用”一个假命题当梯子，用完再把它拆掉。

\begin{toolbox}
本附录的证明语法速查：
\begin{itemize}
\item 命题：能判真假的句子；逆命题是条件与结论对调，真假要独立判断。
\item $p\Rightarrow q$：$p$ 充分，$q$ 必要；两者互为充要时写 $p\Leftrightarrow q$。
\item $\forall$ 的否定是 $\exists$“不”，$\exists$ 的否定是 $\forall$“不”。
\item 推翻“所有”只需一个反例；证明“所有”可以用归纳法（起点加传递）。
\item 正面难攻时用反证法：假设结论不成立，推出矛盾。
\end{itemize}
\end{toolbox}

\section*{思考题}

\textbf{观察题}
\begin{sikao}
\item 写出命题“正方形的四条边都相等”的逆命题，并判断它的真假。提示：找一个四条边都相等但不是正方形的图形。
\end{sikao}

\textbf{证明题}
\begin{sikao}
\item 用数学归纳法证明 $1+2+3+\cdots+n=\dfrac{n(n+1)}{2}$。提示：传递那一步只需在两边同时加上 $k+1$，再通分整理。
\end{sikao}

\textbf{挑战题}
\begin{sikao}
\item 用反证法证明：$\sqrt{2}$ 不能写成两个整数的比。提示：假设 $\sqrt{2}=\dfrac{p}{q}$ 且分数已约到最简，两边平方后考察 $p$ 的奇偶性。
\end{sikao}
